% Generated draft: requires mathematical and source review.
\documentclass[11pt]{article}
\usepackage[margin=25mm]{geometry}
\usepackage{fontspec}
\setmainfont{Latin Modern Roman}
\usepackage{amsmath,amssymb,mathtools,unicode-math}
\setmathfont{Latin Modern Math}
\usepackage{enumitem,needspace,xurl,hyperref}
\setlength{\parindent}{0pt}
\setlength{\parskip}{5pt}
\setlength{\emergencystretch}{3em}
\newcommand{\N}{\mathbb{N}}
\newcommand{\Z}{\mathbb{Z}}
\newcommand{\Q}{\mathbb{Q}}
\newcommand{\R}{\mathbb{R}}
\newcommand{\C}{\mathbb{C}}
\newcommand{\F}{\mathbb{F}}
\newcommand{\A}{\mathbb{A}}
\newcommand{\PP}{\mathbb{P}}
\newcommand{\E}{\mathbb{E}}
\newcommand{\eps}{\varepsilon}
\newcommand{\dd}{\,\mathrm{d}}
\begin{document}
\section*{UnsolvedMath ID: 10400060}
\textbf{Source problem code:} catalogue entry
\subsection{Definitions and mathematical statement}
A virtual knot is conventionally represented by a diagram with classical and virtual crossings, modulo the corresponding Reidemeister and detour moves. The source asks for a series of configuration-space integrals taking values in the arrow-diagram space written in the extraction as \(\overrightarrow{\mathcal A}(I)\), and giving all finite-type invariants of virtual knots.

After the intended diagram relations and finite-type filtration are fixed, a standard universal-series formulation would construct
\[
 Z_{\mathrm{virt}}(K)=\sum_{d\ge0}Z_d(K)
 \quad\text{in}\quad
 \prod_{d\ge0}\overrightarrow{\mathcal A}_d(I),
\]
with coefficients defined by configuration-space integrals, invariant under the allowed virtual-knot moves. Universality means that each finite-type invariant \(v\) of degree at most \(d\) is recovered by a suitable linear functional from the truncated series \(Z_{\le d}\), in the chosen coefficient convention. Merely producing some invariants or a formal arrow-diagram sum would not fulfill that integral/universality request.

The excerpt does not define \(\overrightarrow{\mathcal A}(I)\), its relations, framing or completion, or whether the interval \(I\) refers to long/based virtual knots rather than closed knots. It also does not identify the virtual finite-type notion: different crossing resolutions and filtrations give different theories. No particular theory or arrow relation is inserted.

The geometric configuration spaces, differential forms, collision compactifications and boundary corrections are part of the requested construction and are not supplied. Their invariance must respect the actual virtual-knot equivalence; an arbitrary choice of a surface or diagram representative is not by itself a definition of an invariant.

The summary calls the question open/triage and not obviously resolved, without a precise modern theorem. The 2002 request and later assessment remain source attributions pending current primary-literature review.
\subsection{Short English statement}
Construct configuration-space integrals forming a universal arrow-diagram-valued invariant for the intended finite-type theory of virtual knots. The virtual-knot category, diagram algebra and finite-type convention must first be specified.
\subsection{Sources}
\begin{itemize}
\item Local registry: \url{data/unsolvedmath/problems.json}, permanent UnsolvedMath ID \(10400060\), source code catalogue entry; supplied statement, background, and research summary, updated 2026-08-06.
\item Ohtsuki, \emph{Problems on invariants of knots and 3-manifolds} (2002), Problem 3.10, attributed to D. Thurston, PDF page 55: \url{https://msp.org/gtm/2002/04/gtm-2002-04-024s.pdf}. Supplied PDF-font-structure extraction accessed 2026-07-29; provenance distinguishes 231 numbered items.
\item No separate definition of the arrow algebra, finite-type convention or modern resolution reference is supplied.
\end{itemize}
\end{document}
